\documentclass{article}
\usepackage{circuitikz}
\usepackage{amsmath}
\usepackage{amssymb}

\begin{document}
\begin{figure}
\begin{center}
\begin{circuitikz}
    \draw (0,4) to [ R, l_=$R_1$] (0,2)
    to[ american voltage source, l_=$5$V ] (0, 0)
    to [ short ] (6, 0)
    to [ short ] (6, 4)
    to [ short ] (5.5, 4);
    \draw (2.5,4) to [short] (2.5,5)
    to [short] (3, 5) 
    to [ closing switch, l_=$s_a$ ] (4, 5)
    to [ R, l_=$R_a$ ] (5.5, 5)
    to [ short ] (5.5, 4);
    to [ short ] (5.5, 4);
    \draw (2.5,4) to [short] (2.5,3)
    to [short] (3, 3) 
    to [ closing switch, l_=$s_b$ ] (4, 3)
    to [ R, l_=$R_b$ ] (5.5, 3)
    to [ short ] (5.5, 4);
    \draw (2.5,4) to [ short ] (-1,4);
    \draw (2,4) to [ R, l_=$R_2$ ] (2,0);
    \draw (0,4) to [short] (-2,4)
    to [open, *-*, l_=$V_{\mathrm{out}}$] (-2,0)
    to [short] (0,0);

\end{circuitikz}
\end{center}
\end{figure}

\section{Output Voltage}
The output voltage of the above circuit has a different value depending on the
state of the switches.
\subsection{$s_a$ open, $s_b$ open}
In this case, there is only one equation to solve:

\begin{align}
    0 & = 5 - IR_1 - IR_2 \notag \\
    I(R_1+R_2) &= 5 \notag \\
    I &= \frac{5}{R_1+R_2} 
\end{align}

Then substituting the result to find $V_\mathrm{out}$:
\begin{align}
V_\mathrm{out} &= IR_2 \notag \\
\therefore V_\mathrm{out} &= \frac{5R_2}{R_1+R_2}
\end{align}

\subsection{$s_a$ closed, $s_b$ open}
Now we have to solve for two loops simultaneously:

\begin{align}
    \left[\begin{array}{c}
        5 \\
        0
    \end{array} \right] & = \left[ \begin{array}{cc}
        R_1+R_2 & -R_2 \\
        -R_2 & R_2 + R_a \\
    \end{array}\right] \left[\begin{array}{c}
        I_1 \\
        I_2
    \end{array} \right] \notag \\
    \frac{1}{R_1R_2 + R_1R_a + R_2R_a}\left[ \begin{array}{cc}
        R_2 + R_a & R_2 \\
        R_2 & R_1+R_2 \\
    \end{array}\right] \left[\begin{array}{c}
        5 \\
        0
    \end{array} \right] & =  \left[\begin{array}{c}
        I_1 \\
        I_2
    \end{array} \right] 
\end{align}

And we are only concerned with $V_\mathrm{out}$, so:
\begin{align}
    V_\mathrm{out} & = R_2(I_1 - I_2) \\
    & = \left[\begin{array}{c c}
    R_2 & -R_2
    \end{array}\right]\left[\begin{array}{c}
        I_1 \\
        I_2
    \end{array} \right] \\
    &= \frac{1}{R_1R_2 + R_1R_a + R_2R_a} \left[\begin{array}{c c}
    R_2 & -R_2
    \end{array}\right]\left[ \begin{array}{cc}
        R_2 + R_a & R_2 \\
        R_2 & R_1+R_2 \\
    \end{array}\right] \left[\begin{array}{c}
        5 \\
        0
    \end{array} \right] \\
    &= \frac{5R_2(R_2+R_a) - 5R_2^2}{R_1R_2 + R_1R_a + R_2R_a} \\
    &= \frac{5R_2R_a}{R_1R_2 + R_1R_a + R_2R_a}
\end{align}

\subsection{$s_a$ open, $s_b$ closed}
Clearly, this is the same as the previous calculation except with $s_b$ closed
and $s_a$ open. Therefore:
$$
V_\mathrm{out} = \frac{5R_2R_b}{R_1R_2 + R_1R_b + R_2R_b}
$$

\section{Operating principle}
Now, ideally we can choose $R_1,R_2,R_a,R_b$ such that we have the desired
output voltages in the specified switch states.

We would like to have:

\begin{center}
\begin{tabular}{cc|c}
    $s_a$ & $s_b$ & $V_\mathrm{vout}$ \\ \hline
    OPEN & OPEN & 4.0 V \\
    CLOSED & OPEN & 2.5 V \\
    OPEN & CLOSED & 1.0 V
\end{tabular}
\end{center}

Optimising for these targets (see \verb|./optimise.py|), and keeping $R_1$ fixed (arbitrarily) at 3.3k$\Omega$,
we get a solution of:
\begin{center}
\begin{tabular}{c|c}
    Component & Value \\ \hline
    $R_1$ & $3.3$ k$\Omega$ \\
    $R_2$ & $13.2$ k$\Omega$ \\
    $R_a$ & $4.4$ k$\Omega$ \\
    $R_b$ & $880$ $\Omega$
\end{tabular}
\end{center}

This was verified in LTSpice.
\end{document}